% COSC 3337 Final Project -- Final Report
% Due Monday, December 7, 11:59 PM. Submit the compiled PDF on Canvas.
% Length: roughly 6--10 pages, excluding appendices.
% Replace every [bracketed placeholder]. Put images in the same folder as this file.

\documentclass[11pt]{article}
\usepackage[margin=1in]{geometry}
\usepackage{graphicx}
\usepackage{enumitem}
\usepackage{tabularx}
\usepackage{booktabs}
\usepackage{xcolor}
\usepackage{listings}
\usepackage[colorlinks=true,linkcolor=blue!50!black,urlcolor=blue!50!black]{hyperref}
\setlist{itemsep=2pt, topsep=4pt}

\lstset{
  language=SQL,
  basicstyle=\ttfamily\small,
  keywordstyle=\bfseries\color{blue!50!black},
  commentstyle=\itshape\color{gray},
  breaklines=true,
  frame=single,
  framesep=4pt,
  columns=fullflexible,
  morekeywords={VIEW, REPLACE, HAVING, EXCEPT, INTERSECT, LEFT, RIGHT, OUTER}
}

% ---------- Fill these in ----------
\newcommand{\appname}{[Application Name]}
\newcommand{\members}{[Member 1, Member 2, Member 3, Member 4]}
\newcommand{\repourl}{https://github.com/[user]/[repo]}
% -----------------------------------

\newcommand{\pk}[1]{\underline{#1}}
\newcommand{\fk}[1]{\textit{#1}$^{\mathrm{FK}}$}
\newcommand{\placeholderfig}[1]{\fbox{\parbox[c][2.4in][c]{0.9\textwidth}{\centering [#1]}}}

\begin{document}

\begin{center}
  {\LARGE\bfseries \appname}\\[4pt]
  {\large COSC 3337 Final Project Report}\\[2pt]
  \members\\[2pt]
  Repository: \url{\repourl}
\end{center}

\section{Overview}
% What the application does, who it is for, and the problem it solves.
[Two or three paragraphs.]

\section{Database Design}

\subsection{Final ER diagram}
\begin{center}
  % \includegraphics[width=\textwidth]{er_diagram_final.png}
  \placeholderfig{Insert final ER diagram}
\end{center}

\subsection{Final relational schema}
\begin{itemize}
  \item \textsc{[Table]}(\pk{[id]}, [attribute], \fk{[other\_id]})
  \item \textsc{[Table]}(\pk{[id]}, [attribute], [attribute])
\end{itemize}

\subsection{Changes since Checkpoint 2}
% What changed in the design, and why. Include changes made in response to feedback.
% If nothing changed, say so.
\begin{itemize}
  \item \textbf{[Change]:} [Why you made it]
\end{itemize}

\section{Query Map}
% One row per required query (at least 8). Across all queries you must use at
% least one outer join, GROUP BY with HAVING, subquery, and set operation.
% The full SQL for each query goes in Appendix A.
\noindent\begin{tabularx}{\textwidth}{@{}rXXl@{}}
\toprule
\# & Question it answers & Screen and user action & SQL features \\
\midrule
Q1 & [Plain-English question] & [e.g., ``My Listings'' page, click ``Show unsold''] & [Outer join] \\
Q2 & [Question] & [Screen / action] & [GROUP BY + HAVING] \\
Q3 & [Question] & [Screen / action] & [Subquery] \\
Q4 & [Question] & [Screen / action] & [Set operation] \\
Q5 & [Question] & [Screen / action] & [ ] \\
Q6 & [Question] & [Screen / action] & [ ] \\
Q7 & [Question] & [Screen / action] & [ ] \\
Q8 & [Question] & [Screen / action] & [ ] \\
\bottomrule
\end{tabularx}

\section{View(s)}
% For each view: what it computes and which part of the application uses it.
\begin{itemize}
  \item \textbf{\texttt{[view\_name]}:} [What it does]. Used by [screen / feature].
\end{itemize}
% The CREATE VIEW statement goes in Appendix B.

\section{Screenshots}
% Show the main screens a user moves through. Caption each one.
\begin{figure}[h]
  \centering
  % \includegraphics[width=0.9\textwidth]{screen_home.png}
  \placeholderfig{Screenshot: home page}
  \caption{[What this screen shows and which queries it runs.]}
\end{figure}

\begin{figure}[h]
  \centering
  % \includegraphics[width=0.9\textwidth]{screen_search.png}
  \placeholderfig{Screenshot: [screen name]}
  \caption{[Caption]}
\end{figure}

\section{Team Contributions}
\noindent\begin{tabularx}{\textwidth}{@{}lX@{}}
\toprule
Member & Contributions \\
\midrule
{[Name]} & [What they did] \\
{[Name]} & [What they did] \\
{[Name]} & [What they did] \\
{[Name]} & [What they did] \\
\bottomrule
\end{tabularx}

\section{Setup Instructions}
% Enough detail for someone else to run your application from the repository.
\begin{enumerate}
  \item \textbf{Requirements:} [e.g., Python 3.12, MySQL 8]
  \item \textbf{Create the database:} [command, e.g., \texttt{mysql -u root -p < schema.sql}]
  \item \textbf{Load the data:} [command, e.g., \texttt{python load\_data.py}]
  \item \textbf{Configure the connection:} [where credentials go, e.g., a \texttt{.env} file; never commit real passwords]
  \item \textbf{Run the application:} [command and the URL to open]
\end{enumerate}

\section*{Extra Credit: Deployed Application (optional)}
% Delete this section if you did not deploy.
Live application: \url{[https://...]}\\
Hosted on: [service]

\section*{Sources and Tools}
% Cite datasets, libraries, frameworks, and tutorials you used.
% Disclose any AI tool use, following the course AI policy.
\begin{itemize}
  \item \textbf{Data:} [Dataset name, URL, license]
  \item \textbf{Libraries and frameworks:} [Names and versions]
  \item \textbf{AI tool use:} [Which tools, and what you used them for; or ``None'']
\end{itemize}

\clearpage
\appendix

\section{SQL for Required Queries}
% Paste the exact SQL your application runs. Use ? or %s where the application
% passes a parameter.

\subsection*{Q1: [Short description]}
\begin{lstlisting}
-- Answers: [plain-English question]
SELECT s.name, i.title
FROM seller s
LEFT JOIN item i ON i.seller_id = s.seller_id
WHERE s.seller_id = ?;
\end{lstlisting}

\subsection*{Q2: [Short description]}
\begin{lstlisting}
-- [SQL]
\end{lstlisting}

% Repeat for Q3--Q8.

\section{View Definitions}
\begin{lstlisting}
CREATE VIEW [view_name] AS
SELECT ...;
\end{lstlisting}

\end{document}
